% Standalone problem document, generated from the adjacent README.md.
% Compile with XeLaTeX. Mathematical content is maintained in Markdown.
\documentclass[11pt,a4paper]{article}
\usepackage[fontsize=10.5pt]{scrextend}
\usepackage[margin=22mm,headheight=15pt,headsep=9mm,footskip=11mm]{geometry}
\usepackage{fontspec}
\setmainfont{texgyrepagella}[Extension=.otf,UprightFont=*-regular,BoldFont=*-bold,ItalicFont=*-italic,BoldItalicFont=*-bolditalic]
\setsansfont{texgyreheros}[Extension=.otf,UprightFont=*-regular,BoldFont=*-bold,ItalicFont=*-italic,BoldItalicFont=*-bolditalic]
\setmonofont{texgyrecursor}[Extension=.otf,UprightFont=*-regular,BoldFont=*-bold,ItalicFont=*-italic,BoldItalicFont=*-bolditalic,Scale=MatchLowercase]
\usepackage{amsmath,amssymb}
\usepackage{unicode-math}
\setmathfont{texgyrepagella-math.otf}
\usepackage{xcolor,microtype,enumitem,needspace,fancyhdr}
\usepackage{bookmark}
\definecolor{ink}{HTML}{17344B}
\definecolor{accent}{HTML}{197D80}
\definecolor{muted}{HTML}{596773}
\hypersetup{colorlinks=true,linkcolor=accent,urlcolor=accent,pdftitle={MF-18 - Imaginary-part
rank in a limiting Green-function matrix
equation},pdfauthor={Open Problems in Numerical Linear Algebra}}
\urlstyle{same}
\setlength{\parindent}{0pt}
\setlength{\parskip}{5pt plus 1pt minus 1pt}
\linespread{1.04}
\setlength{\emergencystretch}{3em}
\widowpenalty=10000
\clubpenalty=10000
\displaywidowpenalty=10000
\allowdisplaybreaks[1]
\setlist{nosep,leftmargin=1.5em}
\providecommand{\tightlist}{\setlength{\itemsep}{0pt}\setlength{\parskip}{0pt}}
\providecommand{\pandocbounded}[1]{#1}
\pagestyle{fancy}
\fancyhf{}
\fancyhead[L]{\sffamily\footnotesize\color{muted}OPEN PROBLEMS IN NUMERICAL LINEAR ALGEBRA}
\fancyhead[R]{\sffamily\bfseries\footnotesize\color{accent}MF-18}
\fancyfoot[L]{\parbox[b]{0.9\textwidth}{\sffamily\fontsize{8}{10}\selectfont\color{muted}Editorial ratings; a bounded search cannot certify that a problem remains unsolved.\\Literature check: 2026-09-08}}
\fancyfoot[R]{\sffamily\footnotesize\color{muted}\thepage}
\renewcommand{\headrulewidth}{0.3pt}
\renewcommand{\headrule}{\hbox to\headwidth{\color{accent}\leaders\hrule height \headrulewidth\hfill}}
\makeatletter
\renewcommand{\section}{\Needspace{5\baselineskip}\@startsection{section}{1}{0pt}{12pt plus 2pt minus 2pt}{4pt}{\sffamily\bfseries\large\color{ink}}}
\renewcommand{\subsection}{\Needspace{4\baselineskip}\@startsection{subsection}{2}{0pt}{9pt plus 1pt minus 1pt}{3pt}{\sffamily\bfseries\normalsize\color{ink}}}
\makeatother
\setcounter{secnumdepth}{0}
\begin{document}
{\sffamily\small\color{muted}Matrix functions and stability\par}
\vspace{3pt}
{\raggedright\hyphenpenalty=10000\exhyphenpenalty=10000\sffamily\bfseries\fontsize{21}{24}\selectfont\color{ink}Imaginary-part
rank in a limiting Green-function matrix equation\par}
\vspace{7pt}
{\sffamily\small\textbf{Difficulty:} challenging\quad\textbf{Importance:} interesting
to specialist\par}
\vspace{4pt}
{\color{accent}\hrule height 0.8pt}
\vspace{6pt}
\textbf{Status:} conjecture; no general resolution located as of
2026-09-08

\subsection{Problem statement}\label{problem-statement}

Let \(n\geq1\), let \(C,D,R,P\in\mathbb C^{n\times n}\), and write
\(M^*\) for conjugate transpose. Suppose \(R=R^*\), \(P=P^*\), and

\[
P+\lambda D^*+\lambda^{-1}D\succ0
\qquad (|\lambda|=1).
\]

For each \(\eta>0\), let \(X_\eta\) be the unique nonsingular solution
of

\[
X_\eta+(C^*+i\eta D^*)X_\eta^{-1}(C+i\eta D)
 =R+i\eta P
\]

satisfying \(\rho(X_\eta^{-1}(C+i\eta D))<1\), where \(\rho\) denotes
spectral radius. Assume the finite limit
\(X_0=\lim_{\eta\downarrow0}X_\eta\) exists and is nonsingular. Assume
the matrix polynomial

\[
\mathcal P_0(\lambda)=\lambda^2C^*-\lambda R+C
\]

is regular, meaning \(\det\mathcal P_0(\lambda)\not\equiv0\). Suppose
its eigenvalues on the unit circle are all simple (algebraic
multiplicity one). Their number is even; denote it by \(2m\).

Is it always true that

\[
\operatorname{rank}\!\left(\frac{X_0-X_0^*}{2i}\right)=m?
\]

\subsection{Why this matters for NLA}\label{why-this-matters-for-nla}

The rank describes the non-Hermitian part of the selected
matrix-equation solution in Green-function computations. It connects a
limiting nonlinear solve to the spectrum of a structured quadratic
polynomial.

The source proves the upper bound by \(m\). Its earlier SIAM paper
proves equality for real \(C,R\), \(P=I\), \(D=0\); the general complex
case remains the question. Existence of the limit is assumed, not
conjectured here.

\subsection{References}\label{references}

\begin{enumerate}
\def\labelenumi{\arabic{enumi}.}
\tightlist
\item
  C.-H. Guo, Y.-C. Kuo, W.-W. Lin,
  \href{https://doi.org/10.1016/j.cam.2012.05.012}{\emph{Numerical
  solution of nonlinear matrix equations arising from Green's function
  calculations in nano research}}, JCAM 236 (2012), 4166--4180: equation
  (1), Theorems 3 and 5, and the conjecture on p.~4172.
  \href{https://uregina.ca/~chguo/JCAM_GuoKuoLin.pdf}{Author
  manuscript}: pp.~1, 5--8.
\item
  The same authors, \href{https://doi.org/10.1137/100814706}{\emph{On a
  nonlinear matrix equation arising in nano research}}, SIMAX 33 (2012),
  235--262, §3: the real-coefficient case.
\end{enumerate}

\subsection{Status check}\label{status-check}

On 2026-09-08, checked both papers' scope and Guo's
\href{https://uregina.ca/~chguo/paper.html}{publication list through
2026}. Searches using the JCAM title, authors, ``Green rank
conjecture,'' ``weakly stabilizing,'' and 2025/2026 found no general
resolution. The source proves only the upper bound in the stated complex
setting. This is a bounded literature check.
\end{document}
